\documentclass{beamer}

\usepackage[utf8]{inputenc}
\input{../helper}

\title{Introdução à Computação \\ hardware, software e sistema operacional}
\author{Alexandre Rademaker}
\date{}

\begin{document}

\begin{frame}
 \maketitle
\end{frame}


\begin{frame}[allowframebreaks]{hardware e software}

  A CPU é a unidade de processamento central, que coordena tudo que é
  executado em um computador.

  Sequencias de comandos armazenados na memória são enviados para a
  CPU para serem executados.

  Mas um computador pode estar rodando vários programas ao mesmo
  tempo!  E todos os programas precisam estar armazenados no
  computador para começar, certo?

  \framebreak

  O \textbf{sistema operacional} é o programa que gerencia o
  computador, controlando como cada software usa cada parte do
  hardware do computador.

  \framebreak

  \begin{columns}
    \begin{column}{0.5\textwidth}
      São exemplos de sistemas operacionais: Linux, Windows, MacOS,
      iOS, Android etc.\vspace{1cm}

      Eles gerenciam todos os outros programas que rodam no computador
      e quais dispositivos eles podem usar. Nele instalamos os
      programas.
    \end{column}
    \begin{column}{0.5\textwidth}
      \includegraphics[width=.8\textwidth]{so.png}      
    \end{column}
  \end{columns}

  \framebreak

  A idéia de vários programas rodando `ao mesmo tempo' é criada pelo
  SO, que permite que cada programa em execução tenha uma fração de
  segundos da CPU por vez.

\end{frame}

\begin{frame}[allowframebreaks]{sistema operacional}

  Os sistemas operacionais modernos são acompanhados de vários
  programas quando instalados.

  Existem os \textbf{utilitários}: gerenciador de arquivos, drivers
  para impressão, gerenciador de configurações, relógio,
  \textbf{interpretador de comandos} etc.

  Existem os \textbf{aplicativos}: editores de texto, editores de
  imagem, navegadores etc.

  Também podemos instalar novos programas desde que compatíveis com o
  SO.

  \framebreak

  Os programas podem ser.

  CLI (Command Line Interface) são aqueles que operam por meio de uma
  interface texto, onde o usuário interage digitando texto e recebendo
  texto como resposta, em um terminal. 

  GUI (Graphical User Interface) são aqueles onde o usuário interage
  por meio de elementos visuais, como botões, menus, janelas e ícones.

  Neste curso vamos aprender a implementar programas CLI.

  \framebreak

  Todo sistema operacional moderno (que roda em laptops e desktops)
  oferece dois utilitários muito importantes: gerenciador de arquivos
  (geralmente GUI) e interpretador de comandos (CLI).

  \framebreak

  O \textbf{gerenciador de arquivos} é usado para manipular os
  arquivos e diretórios armazenados no \textbf{sistema de arquivos}
  mantido pelo SO e armazenado no \textbf{hard-drive} do computador.

  Exemplos: Windows Explorer, Finder do MacOS, Nautilus (Linux).

  Diferente do que ocorre no iOS? Como?

  \framebreak

  O \textbf{interpretador de comandos}, ou shell, permite ao usuário
  interagir com o sistema operacional por meio de uma interface de
  linha de comando. Ele é o responsável por ler os comandos digitados
  pelo usuário e executá-los, além de fornecer informações sobre os
  resultados.

  O shell também gerencia as \textbf{variáveis de ambiente}, que são
  valores que podem ser acessados pelos programas em execução. Além
  disso, ele pode executar \textbf{scripts}, que são arquivos com uma
  sequência de comandos que podem ser executados em lote.

  Exemplos: CMD (Windows), bash, zsh, etc.

\end{frame}

\end{document}

%%% Local Variables:
%%% mode: latex
%%% TeX-master: t
%%% End:
